% =========================
% Resume — Centered Header + Extra Section Spacing
% =========================
\RequirePackage{fix-cm}
\documentclass[10pt]{article}

\usepackage[hmargin=0.75in,vmargin=0.5in]{geometry}
\usepackage[hidelinks]{hyperref}
\usepackage{enumitem}
\usepackage{tabularx}
\usepackage{fontawesome5}
\usepackage{titlesec}
\usepackage{xcolor}
\usepackage{setspace}

% Reduce all used text sizes by 3% to keep the resume on one page.
\renewcommand{\normalsize}{\fontsize{9.7pt}{11.64pt}\selectfont}
\renewcommand{\large}{\fontsize{11.64pt}{13.58pt}\selectfont}
\renewcommand{\LARGE}{\fontsize{16.76pt}{21.34pt}\selectfont}

\pagenumbering{gobble}
\setlength{\parindent}{0pt}
\setlength{\parskip}{0pt}

% Tight bullets
\setlist[itemize]{leftmargin=*, nosep}

% Section formatting + MORE SPACE
\titleformat{\section}{\bfseries\large}{}{0pt}{}[\titlerule]
\titlespacing*{\section}{0pt}{12pt}{5pt}

% Helpers
\newcommand{\contactsep}{\hspace{10pt}|\hspace{10pt}}
\newcommand{\roleline}[2]{\textbf{#1}\hfill \textbf{#2}}
\newcommand{\experienceline}[2]{\textbf{#1}\hfill \textbf{#2}}
\newcommand{\orgline}[2]{\textit{#1}\hfill \textit{#2}}

\begin{document}

% ===== Centered Header =====
\begin{center}
    {\LARGE \textbf{STEVEN YANG}} \\[6pt]
    \faPhone\ (647)469-6927
    \contactsep
    \faEnvelope\ \href{mailto:stevenyang.142857@gmail.com}{stevenyang.142857@gmail.com}
    \contactsep
    \faGlobe\ \href{https://Sytion06.github.io}{Website}
    \contactsep
    \faLinkedin\ \href{https://www.linkedin.com/in/steven-yang-7760293a3/}{linkedin}
\end{center}

% ===== Education =====
\section*{Education}
\roleline{University of Toronto}{Sep. 2024 -- May 2028 (Expected)}
\begin{itemize}
  \item \textit{Honours Bachelor of Science}, Data Science Specialist, Computer Science Major
  \item cGPA: 4.0/4.0
  \item \textbf{Awards:} Dean’s List Scholar (2025, 2026) and Woodsworth College Scholarship Fund (2026)
  \item Relevant Courseworks: Software Design, Software Tools and System Programming, Data Structures and Analysis, Data Science, Probability and Statistics, Linear Algebra, Multivariable Calculus
\end{itemize}

% ===== Research Experience =====
\section*{Research}
\experienceline{Student Researcher}{Sep. 2026 -- Present}\\[-1pt]
\orgline{University of Toronto}{}
\begin{itemize}
  \item Build \textbf{AgentLangBench} to evaluate code generation, formal verification, and runtime performance across programming languages for AI agents.
  \item Surveyed 42 languages and tools for AI agents. Used pinned toolchains and repeatable Linux tests to identify 16 entries that passed all applicable checks without workarounds.
  \item Audited AI-generated solutions across languages to assess whether their contracts captured the original task requirements.
\end{itemize}

% ===== Experience =====
\section*{Experience}

\experienceline{Research and Development Intern}{May 2026 -- August 2026}\\[-1pt]
\orgline{NSFOCUS}{Beijing, China}
\begin{itemize}
  \item Built a Python-based \textbf{cyber-monitoring} pipeline integrating \textbf{Tavily} and \textbf{FOFA} APIs to discover, filter, deduplicate, and risk-score suspicious domains and social-media accounts for security investigation.
  \item Prototyped a client-profile-driven \textbf{threat intelligence} workflow for structured event matching, deduplication, and delivery decision using \textbf{Deepseek LLM}.
  \item Analyzed \textbf{AI Coding Agent}'s architecture and workflow in a research report.
\end{itemize}

% ===== Projects =====
\section*{Projects}

\roleline{NHANES Supplement Use Analysis}{2026}
\orgline{Python, pandas, NumPy, matplotlib, scikit-learn}{}
\begin{itemize}
  \item Analyzed NHANES 2021–2023 health and nutrition datasets in a team of 2 to investigate relationships between demographic, dietary, and supplement-use patterns
  \item Performed large-scale \textbf{data cleaning}, \textbf{preprocessing}, \textbf{feature engineering} across multiple linked survey datasets.
  \item Modeled supplement use with statistical and \textbf{machine learning} methods and visualized predictive factors and model performance.
\end{itemize}

\vspace{6pt}

\roleline{Question Bank Platform}{2026}
\orgline{Java, OpenAI API, Full-stack}{}
\begin{itemize}
  \item Built a full-stack Java application to extract and structure math exam questions from PDFs into a searchable digital database.
  \item Integrated \textbf{OCR} and multimodal \textbf{OpenAI Vision models} to parse question text, diagrams, metadata, and page-level images..
  \item Designed async processing with confidence scoring, manual review flags, PDF segmentation, and structured JSON storage.
\end{itemize}

\vspace{6pt}

\roleline{GroupFlow}{2025}
\orgline{Java, Swing, MongoDB, JUnit}{}
\begin{itemize}
  \item Designed and implemented a collaborative task and group management application using \textbf{Clean Architecture}, separating entities, use cases, and UI layers.
  \item Built core features including group creation, role-based permissions, task assignment, and progress tracking.
  \item Improved code maintainability and testability through interface-driven design and unit testing with \textbf{JUnit}.
  \item Collaborated in a team environment using \textbf{Git-based workflows}, pull requests, and code reviews.
\end{itemize}

% ===== Technical Skills =====
\section*{Technical Skills}
\begin{itemize}
  \item \textbf{Programming Languages:} Python, Java, C, JavaScript, HTML/CSS, Rust
  \item \textbf{Developer Tools:} Codex, Claude Code, Git/GitHub, Linux, VS Code, IntelliJ, Jupyter, LaTeX, PostgreSQL, Docker, Verus, Dafny
  \item \textbf{Frameworks/Libraries:} NumPy, pandas, matplotlib, scikit-learn, SciPy, statsmodels, TensorFlow
\end{itemize}

\end{document}
